%==============================================================================
% Chapter 1 -- Introduction
%==============================================================================

\chapter{Introduction}
\label{ch:introduction}

% -----------------------------------------------------------------------------
% This chapter is deliberately broad: it motivates protein nanoparticle vaccines
% as a vaccine platform, and then traces protein design from its physics-based
% origins to the current machine-learning revolution. The two threads meet in
% Chapter 2, where we ask how modern generative models can design these vaccines
% end to end.
% -----------------------------------------------------------------------------
\section{Preliminaries}
\label{sec:preliminaries}

This section introduces the vocabulary that the remainder of the report assumes,
drawn first from structural biology and then from machine learning.

\subsection{Protein structure}
\label{subsec:protein-structure}

A protein is a linear polymer of amino acids. Every residue contributes the same
backbone unit (N, C$\alpha$, C) and one of twenty side chains, so proteins differ
from one another in the order of those side chains along the chain, known as the
\emph{sequence}. In a given environment the sequence largely determines the
three-dimensional structure that the chain adopts
\parencite{anfinsen_principles_1973}. That structure is conventionally described
at four levels (Figure~\ref{fig:protein-structure}). The sequence itself is the
\emph{primary} structure. Hydrogen-bonding patterns between nearby backbone atoms
give two \emph{secondary} structure motifs, namely $\alpha$-helices and
$\beta$-strands. The packing of these elements against one another gives the
\emph{tertiary} structure, the fold of a single chain. Many proteins become
functional only once several chains associate through non-covalent interfaces,
forming a \emph{quaternary} assembly.

\begin{figure}[t]
  \centering
  \includegraphics[width=0.9\linewidth]{Protein_intro_figure_high_res.png}
  \caption{\textbf{The four levels of protein structure.} The amino-acid sequence
    (primary structure) folds locally into elements such as $\alpha$-helices
    (secondary structure), which pack into the fold of a single polypeptide chain
    (tertiary structure); several chains then associate through non-covalent
    interfaces into an assembly of subunits (quaternary structure). Adapted from
    \textcite{rashid_protein_2015}.}
  \label{fig:protein-structure}
\end{figure}

The quaternary level is the one this work addresses. A symmetric assembly is
generated by applying the rotations of a point group to a repeating unit called
the \emph{asymmetric subunit} (ASU). The point groups that close into a finite
assembly are the cyclic groups $C_n$, the dihedral groups $D_n$ and the three
polyhedral groups: tetrahedral $T$, octahedral $O$ and icosahedral $I$, of order
$12$, $24$ and $60$ respectively (Figure~\ref{fig:symmetries}). The order of the
group fixes how many copies of the ASU the particle contains. Helical symmetry,
 is the one case that does not close: it repeats the
subunit along a screw axis to give an open filament rather than a bounded
particle. For all these classes of symmetric proteins, the chains within an 
ASU need not be identical. Where the ASU is a single chain the assembly is a symmetric
homo-oligomer, but placing two or more distinct chains in the ASU gives a
multi-component particle.

\begin{figure}[t]
  \centering
  \includegraphics[width=0.92\linewidth]{protein_symmetries.png}
  \caption{\textbf{Symmetry in protein assemblies.} Representative structures
    with their rotation axes drawn.
    \textbf{(a)} Cyclic $C_3$: human PCNA (PDB~3JA9).
    \textbf{(b)} Dihedral $D_3$: \emph{P.~falciparum} purine nucleoside
    phosphorylase (PDB~1NW4).
    \textbf{(c)} Tetrahedral $T$: \emph{H.~salinarum} dodecin (PDB~1MOG).
    \textbf{(d)} Octahedral $O$: phosphopantetheine adenylyltransferase
    (PDB~6CHL).
    \textbf{(e)} Icosahedral $I$: satellite tobacco necrosis virus (PDB~1STM).
    \textbf{(f)} Helical: Rad51 filament (PDB~6DJO).
    Adapted from \textcite{rcsb_symmetry}.}
  \label{fig:symmetries}
\end{figure}

\subsection{Transformers and attention}
\label{subsec:transformers}

A transformer \parencite{vaswani_attention_2017} is a neural network composed of
repeated blocks. It operates on a set of $n$ \emph{tokens}, each carrying a
feature vector of width $d$. To give an example, in the models used in this
report a token corresponds to a residue or an atom.

The original architecture stacks two kinds of block. An encoder block holds a
self-attention sub-layer followed by a position-wise feed-forward sub-layer. A
decoder block adds a third sub-layer, which attends over the encoder output, and
masks its own self-attention so that a position cannot attend to later ones.
Every sub-layer in either kind is wrapped in the same way: its output is added
back to its input through a residual connection, and the sum is layer-normalised.
Architectures that followed, including the protein design model used in this report,
have expanded the same basic structure of the encoder block here described.

\begin{figure}[t]
  \centering
  \includegraphics[width=0.55\linewidth]{transformer_architecture.png}
  \caption{\textbf{The original transformer architecture.} The encoder stack is
    on the left and the decoder stack on the right, each repeated $N$ times. An
    encoder block pairs multi-head self-attention with a feed-forward sub-layer;
    a decoder block masks its self-attention and inserts a second attention
    sub-layer that reads the encoder output. Each sub-layer is followed by an
    \textsf{Add \& Norm} step, which is the residual connection and the layer
    normalisation. The embeddings, the positional encodings and the final linear
    and softmax head belong to the sequence-to-sequence task of the original work
    rather than to the block itself. Reproduced from
    \textcite{vaswani_attention_2017}.}
  \label{fig:transformer}
\end{figure}

A block therefore updates each token with a revised feature vector. The attention 
mechanism allows each token to consider the context provided by all other tokens 
in the sequence. The update applied to one
token is a learned, weighted combination of contributions from every other token,
so that a residue can come to be represented in terms of the residues it sits
against rather than by its own identity alone. The feed-forward sub-layer, in
contrast, is applied to each token separately and identically, and the residual
connection and the normalisation likewise act on one token at a time; attention
is therefore the only operation in the block through which tokens exchange
information.

Figure~\ref{fig:attention} sets out the mechanism for a single head. After collecting the
input tokens as the rows of $X \in \mathbb{R}^{n \times d}$, three learned
projections turn each token into a query, a key and a value,
%
\[
  Q = X W_Q^{\top}, \qquad K = X W_K^{\top}, \qquad V = X W_V^{\top},
\]
%
with $Q, K \in \mathbb{R}^{n \times d_k}$ and $V \in \mathbb{R}^{n \times
d_v}$. Every query is scored against every key, and a row-wise softmax turns
those scores into the attention matrix
%
\[
  A = \mathrm{softmax}\!\left(\frac{QK^{\top}}{\sqrt{d_k}}\right)
  \in \mathbb{R}^{n \times n},
\]
%
whose $i$th row holds the weights that token $i$ assigns to each token in the
set. Dividing by $\sqrt{d_k}$ keeps the scores in a range where the softmax does
not saturate. The output of the head is $Z = AV$: every token is replaced by a
weighted average of the values, the weights being those learned scores. Several
such \emph{heads} can operate in parallel on different projections of the same
input (not shown in the image); their outputs are concatenated and mapped back to width $d$ by a further
learned projection, and the result enters the residual connection and the
normalisation described above.

\begin{figure}[t]
  \centering
  \includegraphics[width=\linewidth]{Attention.png}
  \caption{\textbf{One attention sub-layer, unrolled.} The multi-head attention
    block at the top is reproduced from the original transformer paper
    \parencite{vaswani_attention_2017}; the panel below expands it for a single
    head and includes the residual connection and the layer normalisation that
    close the sub-layer. Tensors are shaded by role and annotated with their
    dimensions; the symbols are those used in the text.}
  \label{fig:attention}
\end{figure}

The meory cost associated with this vanilla implementation is quadratic in the sequence length. 
Since every token is scored against every other, the scores form an $n\times n$ matrix, so an
attention layer requires $\mathcal{O}(n^2)$ time, and materialising that matrix
requires $\mathcal{O}(n^2)$ memory. On top of this, models that predict or generate the
structure of proteins often introduce a second quadratic object, a \emph{pair representation} holding one
feature vector for every pair of tokens. The attention matrix is formed and
discarded within a single layer, whereas the pair representation is maintained and
updated throughout the network and must therefore be held in memory for the entire
forward pass. We describe in more details this issue and our solution to it in Chapter~\ref{ch:methods}.
%These two quadratic terms are the reason large assemblies are
%expensive to model, and Chapter~\ref{ch:methods} distributes them across several
%GPUs.

\section{Protein nanoparticle vaccines}
\label{sec:pnp-vaccines}

The clinical success of mRNA vaccines during the SARS-CoV-2 pandemic
\parencite{polack_safety_2020} established nucleic-acid delivery as a fast,
programmable vaccine modality: manufacture is sequence-independent and scalable
\parencite{pardi_mrna_2018}, and clinical-grade material for mRNA-1273 entered
production five days after the SARS-CoV-2 sequence was released, reaching a
phase~1 trial 66 days later \parencite{corbett_sars-cov-2_2020}. This platform 
has, however, innately limited immunogenicity. An antigen presented as a
soluble monomer is often only weakly immunogenic, because the size, geometry and
repetitiveness with which an antigen is displayed largely govern how strongly it
engages the immune system \parencite{bachmann_vaccine_2010}.

A complementary strategy is therefore to present the antigen not as a soluble
monomer but arrayed on the surface of a self-assembling \emph{protein
nanoparticle}, in many copies and in a dense, repetitive, virus-like geometry
(Figure~\ref{fig:pnp-vaccine}).
Designed nanoparticle vaccines of this kind have elicited neutralising-antibody
responses roughly an order of magnitude stronger than the corresponding soluble
antigen, both for SARS-CoV-2
\parencite{walls_elicitation_2020} and for respiratory syncytial virus
\parencite{marcandalliInductionPotentNeutralizing2019a}; in the SARS-CoV-2 case
the titres were still undiminished twenty weeks after the boost. The same
principle of geometry-matched multivalent display also underpins designed
miniprotein inhibitors, which neutralise variants of concern and protect mice
when given intranasally before or after challenge
\parencite{hunt_multivalent_2022}. Natural and engineered self-assembling
scaffolds such as lumazine synthase provide versatile platforms for multivalent,
cross-protective antigen presentation
\parencite{ladenstein_second_2020,joseph_lumazine_2025}, and self-assembling
protein nanoparticles are now engineered both for therapeutic delivery
\parencite{olshefskyEngineeringSelfAssemblingProtein2022} and as modular nanocages that
display antibodies and Fc fusions at controlled valency, a format that enhances
receptor agonism and viral neutralisation \parencite{divine_designed_2021}.\\

\begin{figure}[t]
  \centering
  \includegraphics[width=0.58\linewidth]{Protein_nanoparticle_vaccine.jpg}
  \caption{\textbf{Antigen display on a designed nanoparticle.} \textbf{Top:} a
    soluble SARS-CoV-2 receptor-binding domain engages few B-cell receptors,
    whereas the same antigen displayed in sixty copies on a nanoparticle engages
    many at once and elicits a far stronger response. \textbf{Bottom:} the
    scaffold is two-component, a trimer bearing the antigen and a pentamer, which
    are expressed separately and co-assemble into the icosahedral particle.
    Reproduced from \textcite{walls_elicitation_2020}.}
  \label{fig:pnp-vaccine}
\end{figure}

Crucially, because the scaffold is itself a protein, the nanoparticle and its
displayed antigen can be genetically encoded and delivered as a single mRNA
construct, so that the immunogen self-assembles \emph{in vivo}. This
``mRNA-launched'' route combines the manufacturing speed of nucleic-acid vaccines
with the superior immunogenicity of particulate display, and computationally
designed mRNA-launched nanoparticle immunogens have been shown to elicit
protective antibody and T-cell responses in mice
\parencite{hendricks_computationally_2025}. The vaccine platform we target in this
work is therefore \emph{fully protein based}: making it a design target that, as we argue in
Chapter~\ref{ch:comp-design}, can be tackled with modern generative protein models.


% Suggested figure: a schematic contrasting (i) soluble antigen, (ii) nanoparticle
% display, and (iii) mRNA-launched self-assembly. Drop the file in figures/ and
% uncomment.
% \begin{figure}[t]
%   \centering
%   \includegraphics[width=\linewidth]{fig_nanoparticle_vaccine_schematic.pdf}
%   \caption{From soluble antigen to mRNA-launched nanoparticle display.}
%   \label{fig:pnp-schematic}
% \end{figure}


\section{Protein design}
\label{sec:protein-design}

Designing the scaffolds described above is a protein-design problem: we must
specify an amino-acid sequence that folds into, and assembles into, 
a prescribed three-dimensional architecture. This section sketches how the field 
of protein design in general has moved from
physics-based methods to today's machine-learning approaches.

\subsection{Physics-based design}
\label{subsec:physics-design}

The first generation of computational protein design was grounded in biophysics:
candidate sequences were threaded onto a target backbone and scored with an
energy function approximating the physics of folding, with the conformational and
sequence search handled by algorithms such as those in the Rosetta suite. This
paradigm produced landmark results, including the first de novo globular fold
designed to atomic-level accuracy \parencite{kuhlman_design_2003}, de novo enzymes
for reactions with no natural catalyst
\parencite{rothlisberger_kemp_2008,siegel_computational_2010} and picomolar
miniprotein inhibitors of SARS-CoV-2 \parencite{cao_novo_2020}.

For assemblies, symmetry was built into the method itself. The earliest approach
fused two natural oligomers through a rigid helical linker, so that the symmetry
of each component fixed the geometry of the whole
\parencite{padilla_nanohedra_2001}. It gave way to the two-stage recipe that has
defined the field since: rigid-body docking of symmetric building blocks into a
target architecture, followed by atomic-level design of the interfaces this
creates between them. The recipe produced self-assembling nanomaterials whose
crystal structures matched the design models
\parencite{king_computational_2012}, and went on to yield co-assembling
multi-component materials \parencite{king_accurate_2014}, megadalton-scale
two-component icosahedra \parencite{bale_accurate_2016} and hyperstable
60-subunit cages \parencite{hsia_design_2016}, with fast sequence-independent
docking software making the pose search practical
\parencite{sheffler_fast_2023}. These methods work, but they are bounded by their building blocks: docking
draws only on oligomers that already exist, each of which must then be
repurposed, with a new interface designed onto a surface that evolved for
something else. %TO

\subsection{The machine-learning revolution in de novo design}
\label{subsec:ml-design}

The accuracy revolution in protein \emph{structure prediction} that includes AlphaFold2
\parencite{jumperHighlyAccurateProtein2021}, RoseTTAFold
\parencite{baekAccuratePredictionProtein2021} and, more recently, the all-atom
biomolecular model AlphaFold3 \parencite{abramsonAccurateStructurePrediction2024} transformed what is computationally tractable, and it was quickly followed by an
analogous revolution in \emph{generative} design. Deep generative models now propose novel structures directly. Denoising-based
design frameworks including RFdiffusion \parencite{watsonNovoDesignProtein2023},
Chroma \parencite{ingrahamIlluminatingProteinSpace2023} and Genie
\parencite{lin_generating_2023,lin_out_2024,linFastUltraCapableProtein2026},
together with the flow-based Proteina \parencite{geffner_proteina_2025}, enable
the de novo creation of proteins with prescribed structural and functional
properties. Most generate a backbone alone, and their sequences are supplied by a
learned inverse-folding model such as ProteinMPNN
\parencite{dauparasRobustDeepLearning2022}; Chroma instead designs sequence and
side chains for whole complexes itself, and the recent RFdiffusion3
\parencite{butcher_novo_2025}, on which this work builds, generates every atom.

The reach of this toolkit is clearest in what it has been used to make: de novo
minibinders against therapeutically relevant targets, from bioactive peptide
hormones \parencite{vazquez_torres_novo_2024} to clostridial toxins
\parencite{ragotte_novo_2025}; epitope-targeted antibodies, from diffusion-based
co-design of complementarity-determining regions on a fixed framework
\parencite{luo_antigen-specific_2022} to atomically accurate designs of VHHs and
scFvs \parencite{bennett_atomically_2026}; and de novo enzymes
\parencite{yeh_novo_2023}. Recent reviews chart the trajectory of the field
\parencite[e.g.][]{yang_past_2026}.

Two threads run through this chapter: a vaccine platform that is entirely protein
based, and a design toolkit that no longer depends on parts borrowed from nature.
Chapter~\ref{ch:comp-design} brings them together and sets out what still stands
between them and a designed nanoparticle vaccine.
